%=======================================================================
% Requires: booktabs, amsmath, amssymb, siunitx
% Bib keys: s3e, cumulti, opv2v, v2v4real, dairv2x, inscope, deepsense6g,
%           kitti, nuscenes, waymo
%=======================================================================
\section{Related Work}
\label{sec:related}

\subsection{Trajectory Design in Multi-Robot Datasets}
Paradigm-driven collection originates in multi-robot SLAM. S3E~\cite{s3e}
names four collaborative trajectory paradigms and grades each against
loop-closure opportunity using filled and open circles, and
CU-Multi~\cite{cumulti} treats inter-trajectory overlap as a controlled design
variable. In both, the geometry is deliberate, which is the property we retain.
What neither supplies is a measurement of the condition the design produced:
S3E's grading is an assertion about a run rather than a quantity attached to a
frame, and CU-Multi's overlap is obtained by re-running a single platform four
times, so the robots are never simultaneously present and no instantaneous
collaboration geometry exists to measure. Both, moreover, collaborate between
sensors at a common height, $\Delta h = 0$, so viewpoint asymmetry between an
elevated and a ground-level vantage is absent by construction.

\subsection{Collaboration Geometry in Collaborative Perception Datasets}
Collaborative perception datasets treat the relative geometry between agents as
an emergent property of the recording. In vehicle-to-vehicle settings it follows
from how the collecting vehicles happened to drive, whether simulated
\cite{opv2v} or recorded on public roads \cite{v2v4real}; in
vehicle-to-infrastructure settings, from sensor placement and the ego vehicle's
approach \cite{dairv2x}. Per-frame poses are necessarily released, since feature
warping depends on them, so the relative placement of the agents is recoverable.
What is absent is the condition itself: no field states whether the ego's view of
a given object was blocked, whether a partner's was clear, or at what angle the
two observed it, and in the recorded sets these cannot be reconstructed after the
fact, no dense scene model accompanying the release. Splits are cut by town, road
type, or scene, never by collaboration geometry. Even where relative pose is
deliberately diversified at collection time \cite{v2v4real}, it is marginalised
at evaluation: reported accuracy averages over whatever mixture of
configurations a recording contained.

Where geometry does enter this literature, it enters as a nuisance to be
corrected rather than a condition to be reported. Pose error is modelled,
injected, and compensated; the geometry that pose error perturbs is not itself an
axis of the evaluation.

\subsection{Visibility and Evidence as Released Attributes}
Single-agent detection benchmarks do release condition attributes, and both
halves of the formulation used here have precedent there. Geometric visibility
appears as per-object occlusion and truncation flags~\cite{kitti} and as binned
visibility attributes~\cite{nuscenes}. Evidence sufficiency appears as a
filtering rule: annotations carry the count of sensor returns falling inside
them, and objects below a threshold are excluded from scoring or assigned to a
lower difficulty tier~\cite{nuscenes,waymo}. These benchmarks, however, compute
both for a single ego and combine them into one difficulty label, which mixes
occlusion, truncation, and a range proxy into a single tier. The distinction
between an object the ego could not see and one it could see but returned too
little evidence for is thereby lost, and it is that distinction, computed for
every node rather than for one, that the present release keeps.

Occlusion has also been quantified on infrastructure sensors: InScope
\cite{inscope} reports an occlusion degree $\xi_D$ over roadside LiDAR. It is
computed for a fixed roadside vantage, between sensors at a common height, and
summarised per dataset rather than released per frame, so it characterises the
dataset but cannot condition an evaluation within it.

\subsection{Joint Sensing and Communication}
Datasets combining sensing with communication measure blockage at scale, but
infer it from received power itself~\cite{deepsense6g}, so the label and the
signal it is meant to explain are not independent. Computing link state
geometrically, as done here, permits the two to be compared: a geometric
predicate derived from the map is validated against the measured RSSI it never
saw (Section~\ref{sec:results}).

\subsection{Position of This Work}
Four properties of a dataset's collaboration geometry are independent: whether it
is designed, whether it is described, whether it is measured per observation, and
whether the measurement is released in a form that can condition an evaluation.
Multi-robot datasets design and describe it; collaborative perception datasets do
neither; infrastructure datasets measure one component and summarise it away;
sensing-and-communication datasets measure the link but derive the label from the
signal. This work designs the geometry, measures it per (frame, node, object),
validates the measurement against independent signals, and releases it as
columns alongside the recordings. The contribution is the axis and its
validation, not scale: the release is smaller than every dataset named above.
